\documentclass[11pt]{article}\usepackage[margin=1in]{geometry}\usepackage{amsmath,amsthm}\newtheorem{theorem}{Theorem}\newtheorem{remark}{Remark}
\title{The Nonlinear Dispersion-Transfer Frontier}\author{}\date{21 July 2026}\begin{document}\maketitle
\begin{abstract}We compose the nonlinear geometry, normalized comparison, repair audit, retained-history interface, and topology results of Stages 41--49.\end{abstract}
\begin{theorem}[Strongest combined statement]On the certified local chart, nonlinear Fr\'echet dispersion admits an exact gauge-invariant event inclusion into the dependent Gaussian quadratic form. Under a normalized retained-state density comparison and chart-complement bound, the resulting continuous-channel row is at most
\[7.757562504510382\times10^{-5},
\]
inside its allocated budget. A complete positive-mass history table would therefore imply the certified adaptive cluster bound.\end{theorem}
\begin{proof}Stages 41--42 prove the geometric inclusion and normalized tail transfer. Stage 32 supplies the low row, and Stage 46 composes both rows with the chart complement before invoking Stage 36.\end{proof}
The direct repair route does not supply the missing premise: lifting a single-star contraction to links has spill coefficient $1.635983\ldots$, whereas the upper-split sufficient threshold is $0.968323\ldots$.
\begin{remark}The boundary-uniform density and complement estimates remain unproved. Center barriers are not sector partition ratios; cut/orbit control and nonabelian integer topology also remain open. Hence no physical one-scale basin, RG trajectory, continuum theory, or mass gap follows.\end{remark}\end{document}
